Una barra idealizada de sección constante ocupa \(0\le x\le L\). Su temperatura en exceso respecto de un mismo ambiente constante es \(u(x,t)\). El flujo de calor por unidad de área, positivo hacia la derecha, y el balance local son
\[
q(x,t)=-k(x)u_x(x,t),\qquad c_0u_t=-q_x,
\qquad k(x)=k_0\left(1+\frac{x}{L}\right),
\]
con \(L,k_0,c_0>0\). En cada extremo, el flujo \emph{saliente} por unidad de área es proporcional al exceso de temperatura: el coeficiente es \(h_0>0\) a la izquierda y \(h_L>0\) a la derecha. Los coeficientes de transferencia tienen las unidades habituales de flujo por unidad de temperatura.

Traduzca estas leyes a las dos condiciones de contorno para \(u\), justificando sus signos con las direcciones salientes. Para una solución separada \(u(x,t)=X(x)e^{-\lambda t}\), formule después el problema para \(X\) en la convención
\[
-(pX')'+rX=\lambda wX.
\]
Identifique \(p,r,w\), compruebe la regularidad y positividad exigidas en todo el intervalo cerrado e indique las unidades de \(\lambda\).
